\documentclass[a4paper,12pt]{article}
\usepackage[utf8]{inputenc}
\usepackage[T1]{fontenc}
\usepackage[ngerman]{babel}
\usepackage{amsmath}
\usepackage{amssymb}
\usepackage{enumitem}
\usepackage{geometry}
\geometry{a4paper, margin=2.5cm}

\title{Einsendeaufgaben -- Lektion 4\\[0.5em]
\large Modul 61111: Mathematische Grundlagen}
\author{}
\date{}

\begin{document}
\maketitle

\section*{Aufgabe 4.4}

Sei $\varepsilon > 0$ beliebig. Nach dem Satz des Eudoxos existiert ein $n_0 \in \mathbb{N}$, sodass
\[
    n_0 > \frac{16}{\varepsilon}.
\]

F\"{u}r alle $n \geq n_0$ gilt:
\[
    \left| \frac{2n - 14}{n + 1} - 2 \right|
    = \left| \frac{2n - 14 - 2n - 2}{n + 1} \right|
    = \left| \frac{-16}{n + 1} \right|
    = \frac{16}{n + 1}
    \leq \frac{16}{n_0 + 1}
    \leq \frac{16}{n_0}
    < \frac{16}{\frac{16}{\varepsilon}}
    = \varepsilon
\]

Also gilt $\displaystyle\lim_{n \to \infty} \frac{2n - 14}{n + 1} = 2$.

Wir setzen $\varepsilon = \dfrac{1}{4}$ und bestimmen $N$:
\[
    N > \frac{16}{\frac{1}{4}} = 64.
\]

Wir w\"{a}hlen $N = 65$ und verifizieren die Konvergenz an:

F\"{u}r alle $n \geq N = 64$ gilt
\[
    \left| \frac{2n - 14}{n + 1} - 2 \right|
    \leq \frac{16}{64} = \frac{1}{4} = \varepsilon.
\]

\end{document}
